

When a fixed shuffle is applied to a deck of cards repeatedly, the deck will eventually return to its original order. 
The number of times the shuffle must be applied before this happens is called the \textit{period} of the deck for that shuffle.
In the present task, we consider the inverse question: rather than starting with a shuffle and finding its period, 
you are told, for each possible period, how many decks of cards have that period. Your task is to construct any 
shuffle consistent with this information.

Formally, there is a hidden permutation $f = (f_1, f_2, \dots, f_N)$, where $N$ is a given positive integer. A permutation is a list of numbers where
each number from $1$ to $N$ occurs exactly once. Let $x = x_1x_2 \dots x_N$ be a binary string. We say that $f(x)$ ($f$ applied to $x$) is the new
binary string $x_{f_1}x_{f_2} \dots x_{f_N}$, and the \textit{period} of $x$ is the smallest positive integer $p$ such that 

$$\underbrace{f(f(\cdots f(x)\cdots))}_{p\ \text{times}} =  x.$$

It can be proven that for any permutation $f$ and any binary string $x$ of the same length, the period of $x$ exists. In other words, 
if you apply $f$ to $x$ enough times, you eventually get back $x$.

You are given the number $N$, and for each possible period $p$ you are given the number of binary strings 
 (out of all $2^N$ possible binary strings) that have period $p$. Your task is to find any permutation $f$ that corresponds to the given information, 
 or conclude that there is no such permutation.

